\chapter{Codifica Aritmetica (\textit{Arithmetic Coding})}
\label{ch:codifica-aritmetica}

\section{Principi Fondamentali e Procedura Converter}
\label{sec:principi-fondamentali-converter}

La \textbf{codifica aritmetica} (\textit{Arithmetic Coding}), introdotta formalmente da Rissanen \cite{rissanen1976} e perfezionata da Witten, Neal e Cleary \cite{witten1987}, costituisce una tecnica fondamentale di compressione statistica priva di perdite. Essa supera radicalmente il vincolo architetturale dei codici a prefisso classici (come l'algoritmo di Huffman \cite{huffman1952}): l'obbligo di associare a ciascun singolo simbolo un numero intero di bit.

Nel modello shannoniano \cite{shannon1948}, la quantit\`a ideale di informazione generata da un simbolo $\sigma$ con probabilit\`a a priori $p(\sigma)$ \`e espressa dalla propria auto-informazione:
\begin{equation}
  \ell^*(\sigma) = \log_2 \frac{1}{p(\sigma)} \text{ bit}
\end{equation}
Tale valore \`e in generale un numero reale frazionario. Poich\'e i compressori a livello di singolo simbolo non possono emettere frazioni di bit, essi incorrono in un arrotondamento discreto all'intero superiore per ciascun carattere emesso. Al contrario, la codifica aritmetica elabora l'intera sequenza di testo $T$ in modo globale, mappando l'intero messaggio in un \textbf{unico sotto-intervallo semiaperto} del dominio reale unitario:
\begin{equation}
  [l, l + s) \subset [0, 1)
\end{equation}
dove $l$ rappresenta l'estremo inferiore (\textit{base} o \textit{lower bound}) ed $s$ la larghezza (\textit{size} o ampiezza) dell'intervallo. Il corpo compresso finale consiste nella rappresentazione binaria frazionaria di un singolo punto numerico interno a tale intervallo.

\FloatBarrier
\subsection{Definizione Formale della Procedura CONVERTER}
\label{subsec:definizione-converter}

La generazione fisica dei bit di output viene delegata a una procedura formale denominata \texttt{CONVERTER}:
\begin{equation}
  \text{CONVERTER}(x, k) \longrightarrow \langle b_1, b_2, \dots, b_k \rangle
\end{equation}
La procedura riceve in ingresso due argomenti:
\begin{enumerate}
  \item Un valore reale $x \in [0, 1)$, selezionato internamente all'intervallo finale $[l, l + s)$ generato dal processo di compressione;
  \item Un intero positivo $k \in \mathbb{N}^+$, che definisce il numero di cifre binarie di precisione richieste.
\end{enumerate}

L'output coincide con le prime $k$ cifre dell'espansione binaria frazionaria di $x$:
\begin{equation}
  x \approx \sum_{j=1}^{k} b_j \cdot 2^{-j} = 0.b_1 b_2 \dots b_k, \quad \text{con } b_j \in \{0, 1\}
\end{equation}

\FloatBarrier
\subsection{Algoritmo per Moltiplicazioni Successive per Due}
\label{subsec:algoritmo-moltiplicazioni-due}

Il calcolo effettivo delle cifre binarie frazionarie $b_1, \dots, b_k$ segue l'algoritmo classico di moltiplicazione ricorsiva per la base 2:
\begin{enumerate}
  \item Si pone $x_1 = x$;
  \item Per ciascun indice di iterazione $j = 1, 2, \dots, k$:
  \begin{itemize}
    \item Si raddoppia il valore corrente:
    \begin{equation}
      y_j = 2 \cdot x_j
    \end{equation}
    \item Il bit $b_j$ \`e dato dalla parte intera inferiore di $y_j$:
    \begin{equation}
      b_j = \lfloor y_j \rfloor \in \{0, 1\}
    \end{equation}
    \item Si aggiorna il residuo frazionario per il passo successivo:
    \begin{equation}
      x_{j+1} = y_j - b_j = \text{frac}(y_j)
    \end{equation}
  \end{itemize}
\end{enumerate}

\FloatBarrier
\section{Lemma sull'Approssimazione Binaria e Troncamento}
\label{sec:lemma-troncamento-binario}

L'emissione di una sequenza finita di $d$ bit corrisponde matematicamente a un'operazione di troncamento della parte frazionaria. Il seguente risultato analitico quantifica rigorosamente l'errore commesso.

\begin{lemma}[Lemma del Troncamento Binario]
\label{lem:troncamento-binario}
Sia $x \in [0, 1)$ un numero reale avente espansione binaria frazionaria infinita:
\begin{equation}
  x = 0.b_1 b_2 \dots b_d b_{d+1} b_{d+2} \dots = \sum_{j=1}^{\infty} b_j \cdot 2^{-j}
\end{equation}
Sia $\hat{x}$ la sua approssimazione troncata alle prime $d$ cifre binarie:
\begin{equation}
  \hat{x} = 0.b_1 b_2 \dots b_d 0 0 0 \dots = \sum_{j=1}^{d} b_j \cdot 2^{-j}
\end{equation}
Allora l'errore assoluto di troncamento \`e superiormente delimitato da:
\begin{equation}
  x - \hat{x} \le 2^{-d}
\end{equation}
\end{lemma}

\begin{proof}
La differenza algebrica tra il valore esatto $x$ e il valore troncato $\hat{x}$ \`e pari alla somma infinita dei pesi posizionali dei bit azzerati a partire dall'indice $d+1$:
\begin{equation}
  x - \hat{x} = \sum_{i=1}^{\infty} b_{d+i} \cdot 2^{-(d+i)}
\end{equation}
Poich\'e ogni cifra binaria $b_{d+i}$ appartiene all'insieme $\{0, 1\}$, si ha $b_{d+i} \le 1$ per ogni $i \ge 1$:
\begin{equation}
  x - \hat{x} \le \sum_{i=1}^{\infty} 1 \cdot 2^{-(d+i)}
\end{equation}
Raccogliendo a fattor comune il termine costante $2^{-d}$, indipendente dall'indice di sommatoria:
\begin{equation}
  x - \hat{x} \le 2^{-d} \sum_{i=1}^{\infty} 2^{-i}
\end{equation}
La sommatoria residua coincide con la serie geometrica di ragione $q = 1/2$, la cui somma converge a 1:
\begin{equation}
  \sum_{i=1}^{\infty} 2^{-i} = \sum_{i=1}^{\infty} \left(\frac{1}{2}\right)^i = \frac{\frac{1}{2}}{1 - \frac{1}{2}} = \frac{\frac{1}{2}}{\frac{1}{2}} = 1
\end{equation}
Sostituendo il valore calcolato nella disuguaglianza si ottiene la tesi:
\begin{equation}
  x - \hat{x} \le 2^{-d} \cdot 1 = 2^{-d}
\end{equation}
\end{proof}

\FloatBarrier
\section{Processo di Compressione per Suddivisione Progressiva}
\label{sec:processo-compressione-suddivisione}

La codifica comprime un messaggio trasformando sequenzialmente ciascun simbolo in una contrazione geometrica dell'intervallo reale.

\FloatBarrier
\subsection{Modello di Sorgente e Dati Operativi}
\label{subsec:modello-sorgente-esempio}

Si consideri a titolo didattico una sorgente discreta definita sull'alfabeto ternario $A = \{a, b, c\}$, associata alla seguente distribuzione di probabilit\`a:
\begin{itemize}
  \item $p(a) = \dfrac{1}{2} = 0.50$
  \item $p(b) = \dfrac{1}{4} = 0.25$
  \item $p(c) = \dfrac{1}{4} = 0.25$
\end{itemize}
La somma delle probabilit\`a soddisfa la condizione di normalizzazione:
\begin{equation}
  \sum_{\sigma \in A} p(\sigma) = \frac{1}{2} + \frac{1}{4} + \frac{1}{4} = 1
\end{equation}
L'obiettivo \`e comprimere il testo:
\begin{equation}
  T = a\,b\,c \quad (\text{lunghezza } m = 3)
\end{equation}

\FloatBarrier
\subsection{Regole Generali di Aggiornamento Ricorsivo}
\label{subsec:regole-aggiornamento-intervallo}

All'istante iniziale (passo $k = 0$), prima di analizzare qualsiasi carattere, l'intervallo operativo coincide con l'intero segmento unitario:
\begin{equation}
  l_0 = 0, \quad S_0 = 1 \implies [l_0, l_0 + S_0) = [0, 1)
\end{equation}
Ad ogni passo $k \ge 1$, leggendo il carattere $T[k] = \sigma \in A$, l'intervallo attivo $[l_{k-1}, l_{k-1} + S_{k-1})$ viene partizionato in sotto-segmenti contigui e disgiunti di dimensione proporzionale alle probabilit\`a dei simboli:
\begin{enumerate}
  \item \textbf{Aggiornamento dell'ampiezza (\textit{size}):}
  \begin{equation}
    S_k = S_{k-1} \cdot p(\sigma)
  \end{equation}
  \item \textbf{Aggiornamento della base (\textit{lower bound}):}
  \begin{equation}
    l_k = l_{k-1} + S_{k-1} \cdot \sum_{\tau < \sigma} p(\tau)
  \end{equation}
  dove la sommatoria cumula le probabilit\`a di tutti i simboli $\tau$ che precedono $\sigma$ nell'ordine alfabetico fissato.
\end{enumerate}

\FloatBarrier
\subsection{Tracciamento Passo-Passo per la Stringa $T = abc$}
\label{subsec:tracciamento-passo-passo-abc}

L'esecuzione dettagliata della procedura di codifica per la sequenza $T = a\,b\,c$ procede attraverso tre stadi successivi (\autoref{fig:suddivisione-intervalli-aritmetica}):

\begin{enumerate}
  \item \textbf{Passo 1: Elaborazione del simbolo $T[1] = a$}\\
  L'intervallo di partenza $[0, 1)$ viene ripartito secondo le probabilit\`a dell'alfabeto:
  \begin{itemize}
    \item Sotto-intervallo per $a$: si estende da $0$ fino a $p(a) = \dfrac{1}{2}$;
    \item Sotto-intervallo per $b$: si estende da $\dfrac{1}{2}$ fino a $\dfrac{1}{2} + p(b) = \dfrac{3}{4}$;
    \item Sotto-intervallo per $c$: si estende da $\dfrac{3}{4}$ fino a $\dfrac{3}{4} + p(c) = 1$.
  \end{itemize}
  Poich\'e il simbolo letto \`e $a$, il codificatore seleziona il primo segmento:
  \begin{equation}
    l_1 = 0, \quad S_1 = 1 \cdot p(a) = \frac{1}{2} \implies [l_1, l_1 + S_1) = \left[0, \frac{1}{2}\right)
  \end{equation}

  \item \textbf{Passo 2: Elaborazione del simbolo $T[2] = b$}\\
  L'intervallo corrente $[0, 1/2)$ (ampiezza $S_1 = 1/2$) viene suddiviso internamente:
  \begin{itemize}
    \item Segmento per $a$: ampiezza $p(a) \cdot S_1 = \dfrac{1}{2} \cdot \dfrac{1}{2} = \dfrac{1}{4}$, copre $\left[0, \dfrac{1}{4}\right)$;
    \item Segmento per $b$: ampiezza $p(b) \cdot S_1 = \dfrac{1}{4} \cdot \dfrac{1}{2} = \dfrac{1}{8}$, copre $\left[\dfrac{1}{4}, \dfrac{1}{4} + \dfrac{1}{8}\right) = \left[\dfrac{1}{4}, \dfrac{3}{8}\right)$;
    \item Segmento per $c$: ampiezza $p(c) \cdot S_1 = \dfrac{1}{4} \cdot \dfrac{1}{2} = \dfrac{1}{8}$, copre $\left[\dfrac{3}{8}, \dfrac{1}{2}\right)$.
  \end{itemize}
  Leggendo $b$, il codificatore seleziona il sotto-intervallo centrale:
  \begin{equation}
    l_2 = \frac{1}{4}, \quad S_2 = S_1 \cdot p(b) = \frac{1}{2} \cdot \frac{1}{4} = \frac{1}{8} \implies [l_2, l_2 + S_2) = \left[\frac{1}{4}, \frac{3}{8}\right)
  \end{equation}

  \item \textbf{Passo 3: Elaborazione del simbolo $T[3] = c$}\\
  Per analizzare chiaramente la suddivisione di $[1/4, 3/8)$, si riducono le frazioni al denominatore comune 32: la base \`e $l_2 = 1/4 = 8/32$ e l'ampiezza \`e $S_2 = 1/8 = 4/32$, per cui l'intervallo \`e $[8/32, 12/32)$:
  \begin{itemize}
    \item Segmento per $a$: ampiezza $p(a) \cdot S_2 = \dfrac{1}{2} \cdot \dfrac{1}{8} = \dfrac{2}{32}$, copre $\left[\dfrac{8}{32}, \dfrac{10}{32}\right)$;
    \item Segmento per $b$: ampiezza $p(b) \cdot S_2 = \dfrac{1}{4} \cdot \dfrac{1}{8} = \dfrac{1}{32}$, copre $\left[\dfrac{10}{32}, \dfrac{11}{32}\right)$;
    \item Segmento per $c$: ampiezza $p(c) \cdot S_2 = \dfrac{1}{4} \cdot \dfrac{1}{8} = \dfrac{1}{32}$, copre $\left[\dfrac{11}{32}, \dfrac{12}{32}\right)$.
  \end{itemize}
  Il terzo simbolo \`e $c$, inducendo la selezione del segmento superiore:
  \begin{equation}
    l_3 = \frac{11}{32}, \quad S_3 = S_2 \cdot p(c) = \frac{1}{8} \cdot \frac{1}{4} = \frac{1}{32} \implies [l_3, l_3 + S_3) = \left[\frac{11}{32}, \frac{12}{32}\right)
  \end{equation}
\end{enumerate}

Al termine dell'elaborazione della stringa $T = abc$, l'intervallo finale risultante \`e:
\begin{equation}
  \text{COMPRESS}(T) = \left[ \frac{11}{32}, \frac{12}{32} \right) \quad \text{con } l = \frac{11}{32} \text{ e } s = \frac{1}{32}
\end{equation}

\begin{figure}[H]
\centering
\begin{tikzpicture}[
    scale=0.84,
    every node/.style={transform shape}
]
% Stile comune colonne intervalli
\tikzset{
    box/.style={draw=black!70, thick, rounded corners=1pt},
    guide/.style={dashed, black!45, thick}
}

% --- LINEE GUIDA DI PROIEZIONE (tracciate sullo sfondo) ---
% Da colonna 0 (zoom su segmento a: [0, 2.4]) verso colonna 1 ([0, 4.8])
\draw[guide] (1.2, 0) -- (3.3, 0);
\draw[guide] (1.2, 2.4) -- (3.3, 4.8);

% Da colonna 1 (zoom su segmento b: [2.4, 3.6]) verso colonna 2 ([0, 4.8])
\draw[guide] (4.5, 2.4) -- (6.6, 0);
\draw[guide] (4.5, 3.6) -- (6.6, 4.8);

% Da colonna 2 (zoom su segmento c: [3.6, 4.8]) verso colonna 3 ([0, 4.8])
\draw[guide] (7.8, 3.6) -- (9.9, 0);
\draw[guide] (7.8, 4.8) -- (9.9, 4.8);

% --- COLONNA 0: Intervallo Iniziale [0, 1) ---
\node[font=\bfseries\small, color=navy!90, align=center] at (0.6, 5.5) {Stato Iniziale\\{\footnotesize $k=0$}};

\draw[box] (0, 0) rectangle (1.2, 4.8);
% Segmento a (0 to 2.4, ampiezza 1/2) -> SELEZIONATO
\fill[darkgreen!25] (0, 0) rectangle (1.2, 2.4);
\node[font=\small\bfseries, color=darkgreen!90!black] at (0.6, 1.2) {$a: \frac{1}{2}$};
% Segmento b (2.4 to 3.6, ampiezza 1/4)
\fill[navy!12] (0, 2.4) rectangle (1.2, 3.6);
\node[font=\small\bfseries, color=navy!90] at (0.6, 3.0) {$b: \frac{1}{4}$};
% Segmento c (3.6 to 4.8, ampiezza 1/4)
\fill[crimson!12] (0, 3.6) rectangle (1.2, 4.8);
\node[font=\small\bfseries, color=crimson!90] at (0.6, 4.2) {$c: \frac{1}{4}$};

% Ticks asse colonna 0
\node[left, font=\scriptsize, fill=white, inner sep=1pt] at (0, 0) {$0$};
\node[left, font=\scriptsize, fill=white, inner sep=1pt] at (0, 2.4) {$\frac{1}{2}$};
\node[left, font=\scriptsize, fill=white, inner sep=1pt] at (0, 3.6) {$\frac{3}{4}$};
\node[left, font=\scriptsize, fill=white, inner sep=1pt] at (0, 4.8) {$1$};

% --- COLONNA 1: Passo 1, Simbolo 'a' [0, 1/2) ---
\node[font=\bfseries\small, color=navy!90, align=center] at (3.9, 5.5) {Passo 1\\{\footnotesize Simbolo $a$}};

\draw[box] (3.3, 0) rectangle (4.5, 4.8);
% Segmento a (0 to 2.4, ampiezza 1/4 globale)
\fill[darkgreen!12] (3.3, 0) rectangle (4.5, 2.4);
\node[font=\small\bfseries, color=darkgreen!90!black] at (3.9, 1.2) {$a: \frac{1}{4}$};
% Segmento b (2.4 to 3.6, ampiezza 1/8 globale) -> SELEZIONATO
\fill[navy!30] (3.3, 2.4) rectangle (4.5, 3.6);
\node[font=\small\bfseries, color=navy!90!black] at (3.9, 3.0) {$b: \frac{1}{8}$};
% Segmento c (3.6 to 4.8, ampiezza 1/8 globale)
\fill[crimson!12] (3.3, 3.6) rectangle (4.5, 4.8);
\node[font=\small\bfseries, color=crimson!90] at (3.9, 4.2) {$c: \frac{1}{8}$};

% Ticks asse colonna 1
\node[left, font=\scriptsize, fill=white, inner sep=1pt] at (3.3, 0) {$0$};
\node[left, font=\scriptsize, fill=white, inner sep=1pt] at (3.3, 2.4) {$\frac{1}{4}$};
\node[left, font=\scriptsize, fill=white, inner sep=1pt] at (3.3, 3.6) {$\frac{3}{8}$};
\node[left, font=\scriptsize, fill=white, inner sep=1pt] at (3.3, 4.8) {$\frac{1}{2}$};

% --- COLONNA 2: Passo 2, Simbolo 'b' [1/4, 3/8) = [8/32, 12/32) ---
\node[font=\bfseries\small, color=navy!90, align=center] at (7.2, 5.5) {Passo 2\\{\footnotesize Simbolo $b$}};

\draw[box] (6.6, 0) rectangle (7.8, 4.8);
% Segmento a (0 to 2.4 -> 8/32 to 10/32, ampiezza 2/32)
\fill[darkgreen!12] (6.6, 0) rectangle (7.8, 2.4);
\node[font=\small\bfseries, color=darkgreen!90!black] at (7.2, 1.2) {$a: \frac{2}{32}$};
% Segmento b (2.4 to 3.6 -> 10/32 to 11/32, ampiezza 1/32)
\fill[navy!12] (6.6, 2.4) rectangle (7.8, 3.6);
\node[font=\small\bfseries, color=navy!90] at (7.2, 3.0) {$b: \frac{1}{32}$};
% Segmento c (3.6 to 4.8 -> 11/32 to 12/32, ampiezza 1/32) -> SELEZIONATO
\fill[crimson!30] (6.6, 3.6) rectangle (7.8, 4.8);
\node[font=\small\bfseries, color=crimson!90!black] at (7.2, 4.2) {$c: \frac{1}{32}$};

% Ticks asse colonna 2
\node[left, font=\scriptsize, fill=white, inner sep=1pt] at (6.6, 0) {$\frac{8}{32}$};
\node[left, font=\scriptsize, fill=white, inner sep=1pt] at (6.6, 2.4) {$\frac{10}{32}$};
\node[left, font=\scriptsize, fill=white, inner sep=1pt] at (6.6, 3.6) {$\frac{11}{32}$};
\node[left, font=\scriptsize, fill=white, inner sep=1pt] at (6.6, 4.8) {$\frac{12}{32}$};

% --- COLONNA 3: Passo 3, Risultato Finale $T=abc$ ---
\node[font=\bfseries\small, color=crimson!90, align=center] at (10.5, 5.5) {Intervallo Finale\\{\footnotesize Testo $T=abc$}};

\draw[box, draw=crimson!85, fill=crimson!10, very thick] (9.9, 0) rectangle (11.1, 4.8);

% Ticks asse colonna 3
\node[left, font=\scriptsize, fill=white, inner sep=1pt] at (9.9, 0) {$\frac{11}{32}$};
\node[left, font=\scriptsize, fill=white, inner sep=1pt] at (9.9, 4.8) {$\frac{12}{32}$};

% Linea del punto medio x = 23/64 e marker
\draw[thick, navy!90] (9.9, 2.4) -- (11.1, 2.4);
\fill[navy!90] (10.5, 2.4) circle (2.5pt);

% Callout ordinato a destra dell'intervallo finale
\draw[->, >=stealth, thick, navy!90] (11.1, 2.4) -- ++(0.5, 0) 
    node[right, align=left, font=\scriptsize] {
        \textbf{\color{navy!90}Punto medio: } $x = \frac{23}{64}$\\[2pt]
        \textbf{\color{crimson!90}Ampiezza: } $s = \frac{1}{32}$\\[2pt]
        \textbf{\color{darkgreen!90!black}Output: } \texttt{010111}
    };

\end{tikzpicture}
\caption{Suddivisione progressiva dell'intervallo reale unitario durante la codifica del testo $T = abc$: a ciascun passo il sotto-intervallo selezionato viene dilatato concettualmente e suddiviso proporzionalmente alle probabilit\`a dei simboli.}
\label{fig:suddivisione-intervalli-aritmetica}
\end{figure}

\FloatBarrier
\section[Generazione dell'Output e Punto Medio]{Generazione dell'Output\\Compresso e Punto Medio}
\label{sec:generazione-output-punto-medio}

Ottenuto l'intervallo compresso $[l, l + s)$, il codificatore deve identificare un valore numerico che lo rappresenti univocamente all'interno del flusso di bit finale.

\FloatBarrier
\subsection{Formula di Emissione dell'Output}
\label{subsec:formula-emissione-output}

La regola generale stabilisce di codificare il \textbf{punto medio} dell'intervallo, impiegando un numero finito $d$ di bit calibrato in funzione dell'ampiezza $s$:
\begin{equation}
  \text{OUTPUT: CONVERTER}\left(l + \frac{s}{2}, \, \left\lceil \log_2 \frac{2}{s} \right\rceil\right)
\end{equation}
L'espressione individua due quantit\`a determinanti:
\begin{enumerate}
  \item \textbf{La coordinata del punto centrale dell'intervallo:}
  \begin{equation}
    x = l + \frac{s}{2}
  \end{equation}
  \item \textbf{Il numero di bit di precisione richiesti:}
  \begin{equation}
    d = \left\lceil \log_2 \frac{2}{s} \right\rceil
  \end{equation}
\end{enumerate}

\FloatBarrier
\subsection{Calcolo Analitico per il Caso di Studio $T = abc$}
\label{subsec:calcolo-caso-studio-abc}

Applicando la regola ai parametri calcolati per $T = abc$ ($l = 11/32$, $s = 1/32$):
\begin{itemize}
  \item \textbf{Determinazione del punto medio $x$:}
  \begin{equation}
    x = l + \frac{s}{2} = \frac{11}{32} + \frac{\frac{1}{32}}{2} = \frac{11}{32} + \frac{1}{64} = \frac{22}{64} + \frac{1}{64} = \mathbf{\frac{23}{64}}
  \end{equation}
  \item \textbf{Determinazione del numero di bit $d$:}
  \begin{equation}
    d = \left\lceil \log_2 \frac{2}{s} \right\rceil = \left\lceil \log_2 \frac{2}{\frac{1}{32}} \right\rceil = \lceil \log_2 64 \rceil = \lceil 6 \rceil = \mathbf{6\text{ bit}}
  \end{equation}
\end{itemize}
La chiamata finale alla procedura di conversione risulta quindi:
\begin{equation}
  \text{CONVERTER}\left(\frac{23}{64}, \, 6\right)
\end{equation}

\FloatBarrier
\subsection{Dimostrazione di Correttezza del Troncamento a $d$ Bit}
\label{subsec:dimostrazione-correttezza-troncamento}

Affinch\'e il decodificatore possa ricostruire il messaggio originario senza ambiguit\`a, il valore troncato a $d$ bit $\hat{x}$ emesso dal convertitore deve cadere \textbf{rigorosamente all'interno dell'intervallo target}:
\begin{equation}
  \hat{x} \in [l, l + s)
\end{equation}

\begin{proof}
Si valuti l'errore massimo di troncamento $2^{-d}$, dove $d = \lceil \log_2(2/s) \rceil$:
\begin{equation}
  2^{-\left\lceil \log_2 \frac{2}{s} \right\rceil} = \frac{1}{2^{\left\lceil \log_2 \frac{2}{s} \right\rceil}} \le \frac{1}{2^{\log_2 \frac{2}{s}}} = \frac{1}{\frac{2}{s}} = \frac{s}{2}
\end{equation}
La disuguaglianza deriva dalla monotonia dell'esponenziale di base 2: poich\'e $\lceil y \rceil \ge y$, passando al reciproco vale $2^{-\lceil y \rceil} \le 2^{-y}$. In virt\`u del \autoref{lem:troncamento-binario}, l'errore di troncamento soddisfa:
\begin{equation}
  x - \hat{x} \le 2^{-d} \le \frac{s}{2}
\end{equation}
Poich\'e $\hat{x}$ \`e ottenuto azzerando cifre binarie, vale sempre $\hat{x} \le x$. Inoltre, $x = l + s/2$ \`e il punto medio esatto dell'intervallo, distante $s/2$ da entrambi gli estremi. Combinando le disuguaglianze:
\begin{equation}
  l = x - \frac{s}{2} \le \hat{x} \le x < l + s
\end{equation}
Ne consegue che:
\begin{equation}
  \hat{x} \in [l, l + s)
\end{equation}
garantendo che qualsiasi sequenza decodificata a partire da $\hat{x}$ corrisponda univocamente al testo originale $T$.
\end{proof}

\FloatBarrier
\subsection{Conversione Binaria Esplicita di $\frac{23}{64}$ in 6 Bit}
\label{subsec:conversione-binaria-esplicita}

La generazione effettiva dei bit pu\`o essere ricavata attraverso due modalit\`a equivalenti:

\paragraph{Metodo 1: Scomposizione per Somma di Potenze di Due}
Si esprime il numeratore $23$ come somma di potenze intere decrescenti di 2:
\begin{equation}
  23 = 16 + 4 + 2 + 1 = 2^4 + 2^2 + 2^1 + 2^0
\end{equation}
Dividendo l'espressione per il denominatore $64 = 2^6$:
\begin{equation}
  \frac{23}{64} = \frac{16}{64} + \frac{4}{64} + \frac{2}{64} + \frac{1}{64} = 2^{-2} + 2^{-4} + 2^{-5} + 2^{-6}
\end{equation}
Mappando ciascun peso frazionario sui 6 bit di output:
\begin{itemize}
  \item Peso $2^{-1} = 1/2$: non presente $\longrightarrow \mathbf{0}$
  \item Peso $2^{-2} = 1/4$: presente ($16/64$) $\longrightarrow \mathbf{1}$
  \item Peso $2^{-3} = 1/8$: non presente $\longrightarrow \mathbf{0}$
  \item Peso $2^{-4} = 1/16$: presente ($4/64$) $\longrightarrow \mathbf{1}$
  \item Peso $2^{-5} = 1/32$: presente ($2/64$) $\longrightarrow \mathbf{1}$
  \item Peso $2^{-6} = 1/64$: presente ($1/64$) $\longrightarrow \mathbf{1}$
\end{itemize}
La rappresentazione binaria frazionaria \`e dunque:
\begin{equation}
  \frac{23}{64} = [0.010111]_2 \implies \mathbf{b = 010111}
\end{equation}

\paragraph{Metodo 2: Esecuzione Passo-Passo di CONVERTER}
L'applicazione formale dell'algoritmo di moltiplicazione per 2 produce la medesima sequenza (\autoref{tab:tracciamento-converter}):

\begin{table}[H]
\centering
\small
\caption{Tracciamento iterativo delle moltiplicazioni per 2 eseguite da \texttt{CONVERTER}($23/64, 6$).}
\label{tab:tracciamento-converter}
\resizebox{\columnwidth}{!}{%
\begin{tabular}{ccccc}
\toprule
\textbf{Passo $j$} & \textbf{Ingresso $x_j$} & \textbf{Raddoppio $y_j = 2 x_j$} & \textbf{Bit Estratto $b_j = \lfloor y_j \rfloor$} & \textbf{Residuo $x_{j+1} = y_j - b_j$} \\
\midrule
1 & $23/64$ & $46/64 < 1$ & $\mathbf{0}$ & $46/64$ \\
2 & $46/64$ & $92/64 = 1 + 28/64$ & $\mathbf{1}$ & $28/64$ \\
3 & $28/64$ & $56/64 < 1$ & $\mathbf{0}$ & $56/64$ \\
4 & $56/64$ & $112/64 = 1 + 48/64$ & $\mathbf{1}$ & $48/64$ \\
5 & $48/64$ & $96/64 = 1 + 32/64$ & $\mathbf{1}$ & $32/64$ \\
6 & $32/64$ & $64/64 = 1 + 0$ & $\mathbf{1}$ & $0$ \\
\bottomrule
\end{tabular}%
}
\end{table}

La parola compressa emessa per il testo $T = abc$ \`e la sequenza binaria a 6 bit:
\begin{equation}
  c(abc) = \mathbf{010111}
\end{equation}

\FloatBarrier
\section[Efficienza della Codifica Aritmetica]{Teorema sull'Efficienza\\della Codifica Aritmetica}
\label{sec:teorema-efficienza-aritmetica}

Il pregio cardine della codifica aritmetica risiede nell'efficienza asintotica: il sovraccosto rispetto all'entropia teorica di Shannon non dipende dalla lunghezza del testo $m$, ma \`e limitato da una costante prefissata.

\begin{theorem}[Efficienza dell'Arithmetic Coding]
\label{thm:efficienza-aritmetica}
Sia $T$ un testo di $m$ caratteri generato da una sorgente stazionaria priva di memoria con entropia $H(S)$. La dimensione totale in bit dell'output compresso generato dalla codifica aritmetica soddisfa rigorosamente:
\begin{equation}
  \text{Dimensione Output} < m \cdot H(S) + 2\text{ bit}
\end{equation}
\end{theorem}

\begin{proof}
La dimostrazione procede analizzando analiticamente la dimensione in bit $L$ emessa dalla procedura di codifica:

\begin{enumerate}
  \item \textbf{Espressione del numero di bit emessi:}\\
  La formula di emissione richiede una precisione pari a:
  \begin{equation}
    L = \left\lceil \log_2 \frac{2}{s} \right\rceil
  \end{equation}

  \item \textbf{Espansione algebrica del logaritmo:}\\
  Sfruttando le propriet\`a fondamentali dei logaritmi ($\log_2(2/s) = \log_2 2 - \log_2 s = 1 - \log_2 s$):
  \begin{equation}
    L = \lceil 1 - \log_2 s \rceil
  \end{equation}

  \item \textbf{Maggiorazione della funzione ceiling:}\\
  Per ogni numero reale $y$, vale la maggiorazione stretta $\lceil y \rceil < y + 1$. Applicandola a $y = 1 - \log_2 s$:
  \begin{equation}
    L < (1 - \log_2 s) + 1 = 2 - \log_2 s
  \end{equation}

  \item \textbf{Scomposizione dell'ampiezza finale $s$:}\\
  L'ampiezza dell'intervallo finale $s$ per una sequenza di $m$ caratteri $T[1 \dots m]$ \`e pari al prodotto delle probabilit\`a dei singoli simboli incontrati:
  \begin{equation}
    s = \prod_{i=1}^{m} p(T[i]) = \prod_{\sigma \in A} p(\sigma)^{\#\text{occ}(\sigma)}
  \end{equation}
  dove $\#\text{occ}(\sigma)$ denota il numero di occorrenze del simbolo $\sigma$ nel testo $T$.

  \item \textbf{Applicazione del logaritmo al prodotto:}\\
  Sostituendo l'espressione di $s$ nella disuguaglianza:
  \begin{equation}
    L < 2 - \log_2 \left( \prod_{\sigma \in A} p(\sigma)^{\#\text{occ}(\sigma)} \right) = 2 - \sum_{\sigma \in A} \log_2 \left( p(\sigma)^{\#\text{occ}(\sigma)} \right)
  \end{equation}
  Estraendo l'esponente dalla funzione logaritmo:
  \begin{equation}
    L < 2 - \sum_{\sigma \in A} \#\text{occ}(\sigma) \log_2 p(\sigma)
  \end{equation}

  \item \textbf{Introduzione dell'entropia empirica:}\\
  Ricordando che $-\log_2 p(\sigma) = \log_2 \dfrac{1}{p(\sigma)}$, e moltiplicando e dividendo la sommatoria per la lunghezza complessiva $m$:
  \begin{equation}
    L < 2 + m \sum_{\sigma \in A} \frac{\#\text{occ}(\sigma)}{m} \log_2 \frac{1}{p(\sigma)}
  \end{equation}
  Sotto l'ipotesi standard in cui le frequenze empiriche relative nel testo approssimano la distribuzione della sorgente ($\frac{\#\text{occ}(\sigma)}{m} \approx p(\sigma)$):
  \begin{equation}
    L < 2 + m \sum_{\sigma \in A} p(\sigma) \log_2 \frac{1}{p(\sigma)}
  \end{equation}

  \item \textbf{Chiusura della dimostrazione:}\\
  Riconoscendo nella sommatoria la definizione classica dell'entropia di Shannon $H(S) = \sum_{\sigma \in A} p(\sigma) \log_2 \frac{1}{p(\sigma)}$ \cite{shannon1948}:
  \begin{equation}
    L < m \cdot H(S) + 2\text{ bit}
  \end{equation}
\end{enumerate}
\end{proof}

\FloatBarrier
\section{Analisi Comparativa: Codifica Aritmetica vs Huffman}
\label{sec:confronto-aritmetica-huffman}

Il confronto sistematico tra l'algoritmo di Huffman e la codifica aritmetica sintetizza l'evoluzione concettuale della compressione dati statistica (\autoref{tab:confronto-huffman-aritmetica} e \autoref{fig:schema-confronto-huffman-aritmetica}).

\begin{table}[H]
\centering
\small
\caption{Quadro comparativo dettagliato tra Codifica di Huffman e Codifica Aritmetica.}
\label{tab:confronto-huffman-aritmetica}
\resizebox{\columnwidth}{!}{%
\begin{tabular}{lll}
\toprule
\textbf{Propriet\`a Analizzata} & \textbf{Codifica di Huffman} \cite{huffman1952} & \textbf{Codifica Aritmetica} \cite{rissanen1976,witten1987} \\
\midrule
\textbf{Principio Operativo} & Mappatura simbolo per simbolo & Mappatura globale dell'intero testo \\
\textbf{Granularit\`a del Codice} & Numero intero di bit per simbolo & Frazioni continue di bit per simbolo \\
\textbf{Limite Superiore Dimensione} & $< m \cdot H(S) + m\text{ bit}$ & $< m \cdot H(S) + 2\text{ bit}$ \\
\textbf{Overhead Massimo Totale} & Fino a $+m\text{ bit}$ ($+1\text{ bit/simbolo}$) & Al pi\`u $+2\text{ bit}$ per l'intero file \\
\textbf{Tasso Medio per Simbolo ($L/m$)} & $< H(S) + 1\text{ bit/simbolo}$ & $< H(S) + \frac{2}{m} \xrightarrow{m \to \infty} H(S)$ \\
\textbf{Sorgenti Sbilanciate ($H \to 0$)} & Prestazioni scarse ($\ge 1\text{ bit/simbolo}$) & Ottime (l'overhead di 2 bit svanisce per grandi $m$) \\
\textbf{Complessit\`a di Calcolo} & Molto bassa (lookup e bitwise) & Media/Alta (aritmetica a precisione arbitraria) \\
\bottomrule
\end{tabular}%
}
\end{table}

\begin{figure}[H]
\centering
\begin{tikzpicture}[
    scale=0.92,
    every node/.style={transform shape}
]
% Linea verticale del Minimo Teorico di Shannon comune (si ferma al limite della barra superiore)
\draw[thick, dashed, navy!80] (2.0, -0.4) -- (2.0, 3.3);
\node[below=5pt, align=center, font=\small] at (2.0, -0.4) {
    \textbf{\color{navy!90}Minimo Shannon}\\[2pt]
    \color{navy!80}$m \cdot H(S)$
};

% --- BARRA SUPERIORE: CODIFICA DI HUFFMAN ---
\draw[thick, crimson!85, fill=crimson!10, rounded corners=3pt] 
    (2.0, 2.2) rectangle (10.5, 3.3);
\node at (6.25, 2.75) [align=center, font=\small, color=crimson!90!black] {
    \textbf{Codifica di Huffman: Intervallo Ammissibile}\\[2pt]
    {\footnotesize\color{crimson!85!black}$m \cdot H(S) \;\le\; \text{Dimensione Output} \;<\; m \cdot H(S) + m$}
};

% Limite destro Huffman con tick verticale
\draw[very thick, crimson!85] (10.5, 2.05) -- (10.5, 3.45);
\node[below=3pt, font=\scriptsize\bfseries, color=crimson!90] at (10.5, 2.05) {
    $m \cdot H(S) + m$
};

% Parentesi graffa overhead Huffman (in alto, libera da qualsiasi interferenza)
\draw[thick, decoration={brace}, decorate, crimson!85] 
    (2.0, 3.55) -- (10.5, 3.55) 
    node[midway, above=4pt, font=\small\bfseries, color=crimson!90] {
        Overhead di discretizzazione di Huffman: al pi\`u $+m\text{ bit}$ ($+1\text{ bit/simbolo}$)
    };

% --- BARRA INFERIORE: CODIFICA ARITMETICA ---
\draw[thick, darkgreen!85, fill=darkgreen!15, rounded corners=3pt] 
    (2.0, 0.6) rectangle (4.2, 1.5);
\node at (3.1, 1.05) [align=center, font=\scriptsize\bfseries, color=darkgreen!90!black] {
    Codifica\\Aritmetica
};

% Limite destro Codifica Aritmetica con tick verticale
\draw[very thick, darkgreen!85] (4.2, 0.45) -- (4.2, 1.65);

% Callout descrittivo a destra della barra della codifica aritmetica
\draw[->, >=stealth, thick, darkgreen!85] (4.2, 1.05) -- ++(0.7, 0) 
    node[right, align=left, font=\scriptsize] {
        \textbf{\color{darkgreen!90!black}Tetto massimo:} $m \cdot H(S) + 2\text{ bit}$\\[2pt]
        \textbf{\color{darkgreen!80!black}Overhead limitato:} $< +2\text{ bit}$ \textit{(costante e indipendente da $m$)}
    };

\end{tikzpicture}
\caption{Confronto grafico dei massimi teorici di dimensione: mentre la codifica di Huffman accumula un overhead di discretizzazione proporzionale a $m$ (fino a $+m\text{ bit}$), la codifica aritmetica confina l'errore a una quantit\`a fissa e costante ($< +2\text{ bit}$), convergendo esattamente all'entropia per $m \to \infty$.}
\label{fig:schema-confronto-huffman-aritmetica}
\end{figure}

\FloatBarrier
\subsection{Dinamica dell'Ampiezza $s$ e Bit Emessi}
\label{subsec:dinamica-ampiezza-s}

La dinamica dell'ampiezza residua $s = \prod_{i=1}^m p(T[i])$ fornisce un'intuizione geometrica immediata sul consumo di bit della codifica aritmetica:
\begin{itemize}
  \item \textbf{Se $s \approx 1$:} tale scenario occorre quando i simboli elaborati possiedono una probabilit\`a prossima a 1 (simboli estremamente frequenti). L'intervallo subisce una contrazione trascurabile, e la formula di emissione richiede una quantit\`a minima di bit (\textit{``few bits will be added to $l + s/2$''});
  \item \textbf{Se $s \to 0$:} tale situazione si manifesta con l'aumentare della lunghezza $m$ della sequenza o all'inclusione di caratteri rari. L'intervallo collassa rapidamente, richiedendo un numero proporzionato di bit per discriminare con esattezza il punto medio all'interno della partizione (\textit{``lots of bits''}).
\end{itemize}

Mentre l'algoritmo di Huffman deve necessariamente sacrificare almeno un bit intero per ogni carattere del messaggio anche quando $p(\sigma) = 0.999$, la codifica aritmetica assegna idealmente frazioni di bit commisurate all'auto-informazione reale, eliminando la ridondanza strutturale della discretizzazione.
